\documentclass[11pt]{article}

\usepackage[margin=1in]{geometry}
\usepackage[T1]{fontenc}
\usepackage[utf8]{inputenc}
\usepackage{mathpazo}
\usepackage{microtype}
\usepackage{amsmath,amssymb,amsthm,mathtools}
\usepackage{booktabs,tabularx,array}
\usepackage{enumitem}
\usepackage{xcolor}
\usepackage{hyperref}
\usepackage[nameinlink,noabbrev]{cleveref}

\definecolor{linkblue}{HTML}{163A5F}
\hypersetup{
  colorlinks=true,
  linkcolor=linkblue,
  citecolor=linkblue,
  urlcolor=linkblue,
  pdftitle={A Four-State Laboratory for Causal Dissonance, Insight, and Strategic Learning},
  pdfauthor={Research note}
}

\newtheorem{assumption}{Assumption}[section]
\newtheorem{definition}[assumption]{Definition}
\newtheorem{lemma}[assumption]{Lemma}
\newtheorem{proposition}[assumption]{Proposition}
\newtheorem{theorem}[assumption]{Theorem}
\newtheorem{corollary}[assumption]{Corollary}
\newtheorem{remark}[assumption]{Remark}

\crefname{assumption}{assumption}{assumptions}
\Crefname{assumption}{Assumption}{Assumptions}
\crefname{definition}{definition}{definitions}
\Crefname{definition}{Definition}{Definitions}
\crefname{lemma}{lemma}{lemmas}
\Crefname{lemma}{Lemma}{Lemmas}
\crefname{proposition}{proposition}{propositions}
\Crefname{proposition}{Proposition}{Propositions}
\crefname{theorem}{theorem}{theorems}
\Crefname{theorem}{Theorem}{Theorems}
\crefname{corollary}{corollary}{corollaries}
\Crefname{corollary}{Corollary}{Corollaries}
\crefname{remark}{remark}{remarks}
\Crefname{remark}{Remark}{Remarks}

\newcommand{\emptysetset}{\varnothing}
\newcommand{\ind}{\mathbf 1}
\newcommand{\Prow}{\mathcal P^{r}}
\newcommand{\Pcol}{\mathcal P^{c}}
\newcommand{\Pdiag}{\mathcal P^{\times}}
\newcommand{\Pfull}{\mathcal P^{\mathrm{full}}}

\title{A Four-State Laboratory for\\
Causal Dissonance, Insight, and Strategic Learning}
\author{Research note}
\date{July 2026}

\begin{document}

\maketitle

\begin{abstract}
This note develops a deterministic two-firm, four-state model of causal
dissonance and experience-triggered awareness change.  Firms initially cannot formulate a
customer-contingent causal theory.  Heterogeneous private beliefs induce
different committed offering rules, and the resulting experience
empties the active subset of their current model space.  A successful inquiry
episode is represented by two distinct products: a minimally restorative
awareness refinement, followed by an associated causal formulation.  The
$2\times2$ customer-state geometry makes
the central difficulty explicit: the same dissonant experience supports four
minimal refinements and two equally simple factor-based explanations.  No
truth-blind inquiry rule can infer the correct factor from those data alone.
Robustly accommodating both observationally equivalent causal worlds requires
complete awareness, while a non-exhaustive technology must accept fallibility.
Because a restorative refinement fits its generating experience by
construction, retrospective fit carries no discriminating information:
representation-relative certainty exceeds selection-adjusted warrant by a
computable gap that pooled experience never narrows and a single fresh
differentiated observation closes.  Expressing the comparison behind that gap
requires complete awareness, so no minimally refined firm can audit its own
formulation from within.
The model also shows that strategic action controls the diagnostic content of
experience.  Shared false formulations induce pooled actions and can become
self-confirming; heterogeneous formulations generate revealing competition,
learning, and asymmetric value capture.  These pathologies are knife-edge
properties of perfectly forgiving demand: with any outside option, pooled
misservice destroys observable value, and consistency together with
subjective optimality forces every recurring state to be correctly
served---payoff-relevant self-confirming traps, permanent warrant gaps, and
zero-leverage locks all require demand to absorb misservice completely.  An
inquiry-participation layer then
prices the question that dissonance raises: a delabeled leverage statistic,
computed from retained episodes without naming any candidate representation,
determines when costly inquiry is subjectively justified.  Pooled histories
carry zero leverage, so the awareness trap survives even cheap and infallible
inquiry, while differentiated defeat concentrates both the evidence and the
incentive to inquire in the losing firm.  The note separates representation
generation, causal formulation, Bayesian assessment, reflective judgment, and strategic decision and
interprets that separation in relation to Lonergan's account of inquiry and
insight.
\end{abstract}

\section{Purpose and status of the example}

The note is a formal laboratory rather than a complete paper model.  Its
purpose is to determine what can be established in the smallest environment
that permits several qualitatively different awareness refinements: rows,
columns, diagonals, singleton exceptions, and finer combinations.  The model is
designed to answer five questions.
\begin{enumerate}[label=(\roman*)]
  \item How can strategic experience reveal that every currently expressible
  causal theory is inadequate?
  \item What restrictions does the experience place on a satisfactory new
  representation?
  \item Why do those restrictions direct inquiry without logically determining
  its answer?
  \item Where does fixed-space Bayesian learning stop, and where do
  representation production and subsequent judgment begin?
  \item How do the firms' resulting strategies affect whether a false
  formulation will subsequently be exposed?
  \item When is costly inquiry subjectively justified, and what represented
  statistic can price a question whose possible answers the firm cannot yet
  formulate?
  \item How much confidence does the generating experience actually license in
  a produced formulation, and which agents can compute that warrant?
  \item Which epistemic pathologies survive when misservice destroys value
  rather than merely reallocating it?
\end{enumerate}

Four modeling distinctions are maintained throughout.  An \emph{awareness
partition} determines which customer distinctions a firm can represent.  A
\emph{causal theory} assigns a customer response to every cell of that
partition.  A \emph{causal formulation} combines a newly available partition
with an expressible theory.  A \emph{strategy} maps represented customer states
to product architectures.  None of these objects is identified with insight
itself; the mathematics represents the output of a successful insight-producing
process.

\section{Notation}

\begin{center}
\renewcommand{\arraystretch}{1.15}
\begin{tabularx}{\textwidth}{@{}>{\raggedright\arraybackslash}p{0.28\textwidth}X@{}}
\toprule
\textbf{Symbol} & \textbf{Meaning} \\
\midrule
$N=\{1,2\}$ & Set of firms. \\
$A=\{0,1\}$ & Set of architectures and firm actions. \\
$X=\{x_{00},x_{01},x_{10},x_{11}\}$ & Analyst's finite customer-state space. \\
$d:X\to A$ & Objectively preferred architecture. \\
$u_i$ & Firm $i$'s objective payoff function. \\
$\beta$ & Demand retention under pooled misservice; $\beta=1$ is the baseline. \\
$\Pi(X)$ & Set of all partitions of $X$. \\
$\mathcal P_i\in\Pi(X)$ & Firm $i$'s awareness partition. \\
$M(\mathcal P_i)$ & Set of causal theories expressible under $\mathcal P_i$. \\
$m\in M(\mathcal P_i)$ & An individual causal theory. \\
$\mu_i\in\Delta(M(\mathcal P_i))$ & Firm $i$'s belief over its expressible theories. \\
$\sigma_i:\mathcal P_i\to A$ & Firm $i$'s product policy. \\
$H$ & Set of analyst-level histories. \\
$h$ & Analyst's record of customer states, actions, and payoffs. \\
$e_i(h;\mathcal P_i)$ & Firm $i$'s deliberative record of retained episode tokens. \\
$D(h)\subseteq X\times A$ & Causal labels revealed by differentiated actions. \\
$\widehat M(\mathcal P,D)$ & Theories under $\mathcal P$ consistent with $D$. \\
$E(D)$ & Conflict-edge set induced by $D$. \\
$\mathcal R(\mathcal P,D)$ & Set of minimally restorative refinements. \\
$\kappa(\mathcal P)$ & Coordinate complexity of a binary partition. \\
$\phi_i:\Pi(X)\times H\to\Pi(X)$ & Analyst-level representation-producing technology. \\
$\mathcal I(\mathcal P,D)$ & Complexity-filtered restorative correspondence. \\
$\overline\mu_i$ & Representation-relative coherence distribution. \\
$\varepsilon_i$ & Firm $i$'s binary inquiry election. \\
$\gamma_i$ & Firm $i$'s inquiry activation cost. \\
$\theta_i$ & Firm $i$'s delabeled inquiry-success weight. \\
$L_i(h)$ & Disagreement set: retained labeled episodes contradicting the committed policy. \\
$B_i(h)$ & Retained-record leverage of inquiry. \\
$\mathcal D$ & Candidate objective laws in a reflective assessment. \\
$\nu$ & Reflective prior over $\mathcal D$. \\
$W(m'\mid h)$ & Selection-adjusted warrant of formulation $m'$. \\
$G(m'\mid h)$ & Warrant gap: coherence probability minus warrant. \\
$\Prow,\Pcol,\Pdiag$ & Row, column, and diagonal partitions. \\
$\mathcal P^{jk}$ & Partition isolating $x_{jk}$ from the other states. \\
$\Pfull$ & Partition of $X$ into four singleton cells. \\
\bottomrule
\end{tabularx}
\end{center}

Capital Roman letters denote sets, lowercase Roman and Greek letters denote
elements, parameters, and functions, calligraphic letters denote sets of sets,
and blackboard-bold notation is reserved for canonical sets.  For any finite
set $Z$, $\Delta(Z)$ denotes the set of probability distributions on $Z$.

\section{Objective environment and timing}

\subsection{Customer states and objective fit}

There are two firms,
\[
N=\{1,2\},
\]
and two available architectures,
\[
A=\{0,1\}.
\]
The analyst's customer-state space is
\[
X=\{x_{00},x_{01},x_{10},x_{11}\}.
\]
It is convenient to display the states as a matrix:
\[
\begin{array}{c|cc}
 & c=0 & c=1\\
\hline
r=0 & x_{00} & x_{01}\\
r=1 & x_{10} & x_{11}.
\end{array}
\]
The two coordinates are analyst-level descriptions of potentially usable
customer distinctions.  Their status in the firms' awareness is specified in
\cref{sec:unawareness}.

The objectively preferred architecture is
\begin{equation}
d:X\longrightarrow A,
\qquad
d(x_{rc})=r.
\label{eq:true-row-law}
\end{equation}
Thus the first coordinate is causally relevant and the second is not:
\[
\begin{array}{c|cc}
 & c=0 & c=1\\
\hline
r=0 & 0 & 0\\
r=1 & 1 & 1.
\end{array}
\]

If the firms offer different architectures, the customer buys from the firm
offering $d(x)$.  If they offer the same architecture, they split the unit
market, even if another architecture would have served the customer better.
Firm $i$'s payoff is
\begin{equation}
u_i(a_i,a_j,x)=
\begin{cases}
\frac12, & a_i=a_j,\\
1, & a_i\neq a_j\text{ and }a_i=d(x),\\
0, & a_i\neq a_j\text{ and }a_j=d(x).
\end{cases}
\label{eq:objective-payoff}
\end{equation}
The payoffs sum to one.  A customer state can equivalently be interpreted as a
unit market segment, which avoids interpreting the one-half outcome as random
tie-breaking for an indivisible customer.

\subsection{Batch commitment}

The opening experience occurs within a product generation.  At the start of a
generation, each firm simultaneously commits to a contingent offering or
configuration rule.  A finite batch of customer states is then served under
those rules.  Firms observe the customer episodes, actions, and payoffs, but
they cannot revise the committed rule until the batch ends.  Initial coarse
awareness makes the rule constant; refined awareness permits customer-
contingent implementation.  Updating and inquiry take place only between
generations.

\begin{assumption}[Diagnostic batch]
\label{ass:batch}
The initial batch consists of the two states $x_{00}$ and $x_{11}$.  Both are
served under the policies chosen before either outcome is observed.
\end{assumption}

The batch commitment is substantive.  With deterministic period-by-period
updating and immediate redesign, the first differentiated observation would
collapse both firms' active sets onto the same constant theory.  They would
then choose the same architecture at the second state, suppressing the causal
signal needed for the focal dissonance.

\section{Awareness and expressible causal theories}
\label{sec:unawareness}

\begin{definition}[Awareness partition]
An awareness partition $\mathcal P_i\in\Pi(X)$ determines the customer
distinctions firm $i$ can use in causal theories and prospective product
policies.  A policy is a function
\[
\sigma_i:\mathcal P_i\longrightarrow A.
\]
The action taken at $x$ is $\sigma_i(C)$, where $C\in\mathcal P_i$ is the cell
containing $x$.  For brevity, $\sigma_i(x)$ denotes this induced statewise
action.
\end{definition}

For any $\mathcal P\in\Pi(X)$, let $C_{\mathcal P}(x)$ denote the unique cell
of $\mathcal P$ containing $x$.

Initially, both firms have the coarse awareness partition
\begin{equation}
\mathcal P_i^0=\{X\}.
\label{eq:coarse-partition}
\end{equation}
A firm with this partition cannot formulate a causal theory, or commit to a
policy, that distinguishes any pair of customer states.

Let $n$ denote the claim that architecture is causally neutral.  For any
partition $\mathcal P$, define
\begin{equation}
M(\mathcal P)
=
\left\{
m:X\longrightarrow\{0,1,n\}:
m(x)=m(x')
\text{ whenever }x,x'\text{ lie in the same cell of }\mathcal P
\right\}.
\label{eq:model-space}
\end{equation}
For $C\in\mathcal P$ and $m\in M(\mathcal P)$, write $m(C)$ for the common
value $m(x)$ on $C$, which is well defined by cell constancy.
The interpretations are
\[
m(x)=a:
\text{ architecture }a\text{ is preferred at }x,
\qquad
m(x)=n:
\text{ architecture is neutral at }x.
\]

The payoff predicted by $m$ is
\begin{equation}
u_i^m(a_i,a_j,x)=
\begin{cases}
\frac12, & a_i=a_j\text{ or }m(x)=n,\\
1, & a_i\neq a_j\text{ and }m(x)=a_i,\\
0, & a_i\neq a_j\text{ and }m(x)=a_j.
\end{cases}
\label{eq:subjective-payoff}
\end{equation}

Under $\mathcal P_i^0$, the only expressible theories are the three constant
functions
\[
M(\mathcal P_i^0)=\{m^0,m^1,m^n\},
\]
where $m^a(x)=a$ for every $x\in X$ and $a\in\{0,1,n\}$.  The true law $d$ is
not in $M(\mathcal P_i^0)$.

The partition is an analyst-level representation of the firm's subjective
state space.  At the level of deliberation and expressible belief, the firm has
the quotient state space $X/\mathcal P_i$, whose elements are the cells of
$\mathcal P_i$; it does not possess the master space $X$ together with an
unobserved choice among finer partitions.  Every pre-insight belief,
likelihood, deliberately entertained question, and policy is
$\mathcal P_i$-measurable.  In particular, before awareness expands, the firm
cannot name a finer partition, assign it a probability, or condition a policy
on it.

\begin{assumption}[Retention and recodability]
\label{ass:recodability}
Customer episodes can be retained as particulars and retrospectively recoded
after the firm acquires a new predicate.  Once acquired, the predicate can also
classify a future customer before the firm acts.  Before acquisition, neither
belief nor action can condition on that predicate.
\end{assumption}

This assumption supplies the minimal distinction between experiencing a
particular and possessing a projectible causal concept about it.  Without some
extra-deliberative access to retained particulars, a strict refinement could
not arise endogenously from a subjective state space that contains only one
state.  Retaining a particular is not possession of an event that separates it
from other customer states and cannot support a contingent policy before
recoding.  The representation-producing technology introduced below is a
primitive of concept production: it can operate on retained particulars through
latent comparison capacities, but those operations are not available to the
firm's pre-insight deliberation.

\section{Private beliefs and the opening product policies}

Firm $i$ has a private prior
\[
\mu_i^0\in\Delta(M(\mathcal P_i^0)).
\]
The priors need not be common knowledge.  For a general awareness cell
$C\in\mathcal P_i$, write
\[
p_i^a(C)
=
\mu_i\big(\{m\in M(\mathcal P_i):m(C)=a\}\big),
\qquad a\in\{0,1,n\}.
\]
Let $v_i(a\mid C)$ denote firm $i$'s subjective expected payoff from action
$a$ at cell $C$, given its current causal belief and any conjecture about the
rival's action.

\begin{assumption}[Subjective optimality and commitment]
\label{ass:subjective-optimality}
Whenever firm $i$ has a well-defined current causal belief, it selects at the
start of each generation an action maximizing $v_i(a\mid C)$ at every
represented cell $C\in\mathcal P_i$ and commits to the resulting policy for the
generation's batch.  The carry-forward clause in
\cref{ass:dissonance-trigger} governs the exceptional state in which
dissonance has made belief undefined and representation production has failed.
\end{assumption}

\begin{lemma}[Subjective dominance]
\label{lem:dominance}
For either pure action of the rival, and therefore for any mixture over the
rival's actions,
\begin{equation}
v_i(0\mid C)-v_i(1\mid C)
=
\frac12\big[p_i^0(C)-p_i^1(C)\big].
\label{eq:dominance-difference}
\end{equation}
Consequently, action $0$ is subjectively strictly dominant on $C$ if
$p_i^0(C)>p_i^1(C)$, and action $1$ is subjectively strictly dominant if the
inequality is reversed.
\end{lemma}

\begin{proof}
Suppose first that the rival selects $0$.  Selecting $0$ yields one-half under
every theory.  Selecting $1$ yields zero under a theory labeling the cell $0$,
one under a theory labeling it $1$, and one-half under a neutral theory.  The
expected-payoff difference is therefore
\[
\frac12-\left[p_i^1(C)+\frac12p_i^n(C)\right]
=
\frac12\big[p_i^0(C)-p_i^1(C)\big].
\]
If the rival selects $1$, action $0$ yields
$p_i^0(C)+\frac12p_i^n(C)$ and action $1$ yields one-half, giving the same
difference.  Linearity establishes the result against any mixture.
\end{proof}

\begin{assumption}[Heterogeneous initial causal beliefs]
\label{ass:heterogeneous-priors}
\[
\mu_1^0(m^0)>\mu_1^0(m^1),
\qquad
\mu_2^0(m^1)>\mu_2^0(m^0).
\]
\end{assumption}

\begin{proposition}[Belief-induced differentiation]
\label{prop:initial-differentiation}
Under \cref{ass:batch,ass:subjective-optimality,ass:heterogeneous-priors}, firm $1$ commits to
architecture $0$ throughout the opening batch and firm $2$ commits to
architecture $1$.
\end{proposition}

\begin{proof}
Under $\mathcal P_i^0$, each product policy must be constant on $X$.
\Cref{lem:dominance} and \cref{ass:heterogeneous-priors} make action $0$
strictly dominant for firm $1$ and action $1$ strictly dominant for firm $2$.
Batch commitment then holds the differentiated action profile fixed while both
diagnostic states are served.
\end{proof}

\section{Revealed causal evidence and Bayesian closure}

Let $H$ be the set of finite analyst-level histories, with a typical history
written
\[
h=\big(x_s,a_{1s},a_{2s},u_{1s},u_{2s}\big)_{s=1}^{t}.
\]
At the deliberative level, firm $i$ represents this history under
$\mathcal P_i$ as
\begin{equation}
e_i(h;\mathcal P_i)
=
\big(s,C_{\mathcal P_i}(x_s),a_{1s},a_{2s},u_{1s},u_{2s}\big)_{s=1}^{t}.
\label{eq:deliberative-record}
\end{equation}
The index $s$ retains an episode as a particular token; it is not a customer
predicate on which a pre-insight theory or policy may condition.  The
analyst-level history also records the underlying customer state so that a
newly produced representation can retrospectively recode the token.

When the actions differ, the payoff vector identifies the winning architecture.
Let
\begin{equation}
D(h)
=
\left\{
(x_s,a_{is}):
s\in\{1,\ldots,t\},\ i\in N,\
a_{1s}\neq a_{2s}\text{ and }u_{is}=1
\right\}
\subseteq X\times A
\label{eq:revealed-data}
\end{equation}
be the analyst's recoding of causal labels attached to the retained episodes.
Before refinement, the firm does not observe the ordered pair $(x_s,a_{is})$ as
a proposition on the master space.  It experiences two episode tokens in its
same represented cell with the displayed actions and opposing payoffs.  After
refinement, \cref{ass:recodability} permits the tokens to be recoded into the
new cells.  A pooled observation is omitted because its payoff vector is
$(1/2,1/2)$ under every objective and subjective causal law.

A revealed dataset is \emph{functional} if
\begin{equation}
(x,a),(x,a')\in D
\quad\Longrightarrow\quad
a=a'.
\label{eq:functional-data}
\end{equation}
Every dataset generated by the deterministic objective law is functional.

For a proposed awareness partition $\mathcal P$, define the active model set
\begin{equation}
\widehat M(\mathcal P,D)
=
\left\{
m\in M(\mathcal P):
m(x)=a\text{ for every }(x,a)\in D
\right\}.
\label{eq:active-set}
\end{equation}

\begin{definition}[Causal dissonance]
An awareness partition $\mathcal P$ exhibits causal dissonance at dataset $D$
when
\[
\widehat M(\mathcal P,D)=\emptysetset.
\]
The expressible model set $M(\mathcal P)$ does not disappear.  Rather, no model
in it accounts for all revealed causal evidence.
\end{definition}

In the opening batch, \cref{eq:true-row-law,prop:initial-differentiation} imply
\begin{equation}
\overline D
=
\{(x_{00},0),(x_{11},1)\}.
\label{eq:discovery-data}
\end{equation}
No constant theory fits both observations, so
\begin{equation}
\widehat M(\mathcal P_i^0,\overline D)=\emptysetset.
\label{eq:initial-dissonance}
\end{equation}
The immediate question has a determinate defect but not a determinate answer:
what difference between $x_{00}$ and $x_{11}$ explains why different
architectures fit them?

\begin{lemma}[Endogenous revelation]
\label{lem:revelation}
At a customer state $x$, the objective label $d(x)$ is revealed by the payoff
vector if the firms choose different architectures.  If they choose the same
architecture, the payoff vector is independent of $d(x)$ and supplies no
restriction on the active model set.
\end{lemma}

\begin{proof}
With two architectures, differentiated actions imply that exactly one firm
offers $d(x)$.  The firm receiving payoff one identifies that architecture.
Under pooled actions, \cref{eq:objective-payoff} gives both firms one-half for
either value of $d(x)$.
\end{proof}

\begin{lemma}[Fixed-space Bayesian closure]
\label{lem:bayesian-closure}
Fix $\mathcal P$ and a prior $\mu\in\Delta(M(\mathcal P))$.  Whenever Bayesian
conditioning has a positive denominator, the posterior remains in
$\Delta(M(\mathcal P))$.  Every theory in posterior support and every optimal
policy therefore remains constant within each cell of $\mathcal P$.  If
$\widehat M(\mathcal P,D)=\emptysetset$ in the deterministic model, the Bayesian
posterior is undefined and conditioning supplies no strict refinement of
$\mathcal P$.
\end{lemma}

\begin{proof}
Conditional on the realized customer-state and action design, define the
deterministic likelihood
\[
\ell(D\mid m,\mathcal P)
=
\ind\{m\in\widehat M(\mathcal P,D)\}.
\]
Whenever the denominator is positive, Bayes' rule gives
\[
\mu(m\mid D,\mathcal P)
=
\frac{\mu(m)\ell(D\mid m,\mathcal P)}
{\sum_{m'\in M(\mathcal P)}
\mu(m')\ell(D\mid m',\mathcal P)}.
\]
Conditioning changes probability weights but not the domain on which those
weights are defined.  By \cref{eq:model-space}, every element of
$M(\mathcal P)$ is constant within $\mathcal P$-cells, so any posterior mixture
and any policy optimized against it respects those cells.  When every model
assigns likelihood zero to the data, Bayes' denominator is zero.  No operation
in the conditioning formula introduces a partition outside the prior domain.
\end{proof}

If the prior lacks full support, a nonempty active set can still receive total
prior mass zero.  The direction used here is the sharper one needed for causal
dissonance: an empty active set necessarily gives a zero denominator under
every prior.

\begin{remark}[Algebraic interpretation]
Let the event algebra generated by $\mathcal P$ consist of all unions of its
cells.  Boolean operations and conditioning within that algebra cannot produce
an event separating two states in the same cell.  A strict awareness refinement
therefore requires a representation-producing primitive that is not reducible
to fixed-space Bayesian conditioning.
\end{remark}

\section{Restorative inquiry and the conflict structure}

For $\mathcal P,\mathcal P'\in\Pi(X)$, write
\[
\mathcal P\prec\mathcal P'
\]
when $\mathcal P'$ strictly refines $\mathcal P$: every cell of
$\mathcal P'$ lies within a cell of $\mathcal P$, and the partitions are not
equal.

\begin{definition}[Restorative refinement]
Partition $\mathcal P'$ is restorative at $(\mathcal P,D)$ if
\[
\mathcal P\prec\mathcal P'
\qquad\text{and}\qquad
\widehat M(\mathcal P',D)\neq\emptysetset.
\]
It is \emph{minimally restorative} if no restorative $\mathcal P''$ satisfies
\[
\mathcal P\prec\mathcal P''\prec\mathcal P'.
\]
Let $\mathcal R(\mathcal P,D)$ denote the set of minimally restorative
refinements.
\end{definition}

Minimality is defined in the refinement partial order.  It does not imply that
there is a unique coarsest answer.

Define the conflict-edge set
\begin{equation}
E(D)
=
\left\{
\{x,x'\}\subseteq X:
(x,0),(x',1)\in D
\text{ or }(x,1),(x',0)\in D
\right\}.
\label{eq:conflict-edges}
\end{equation}
A subset of $X$ is conflict-free if it contains neither endpoint pair of any
edge in $E(D)$.

\begin{proposition}[Conflict characterization]
\label{prop:conflict-characterization}
For any $\mathcal P\in\Pi(X)$ and functional revealed dataset $D$,
\[
\widehat M(\mathcal P,D)\neq\emptysetset
\]
if and only if every cell of $\mathcal P$ is conflict-free.  Consequently,
$\mathcal P'$ is restorative at $(\mathcal P,D)$ if and only if it strictly
refines $\mathcal P$ and every cell of $\mathcal P'$ is conflict-free.
\end{proposition}

\begin{proof}
If a cell contains the endpoints of a conflict edge, the observations require
a cell-constant theory to assign both zero and one to that cell, which is
impossible.  Conversely, suppose every cell is conflict-free.  If a cell
contains a revealed state, all revealed states in it have the same label;
assign that label to the cell.  Assign any value in $\{0,1,n\}$ to an unlabeled
cell.  The resulting cell-constant function belongs to
$\widehat M(\mathcal P,D)$.
\end{proof}

The conflict characterization generalizes beyond the focal example.  The
restorative correspondence identifies minimally refined partitions whose cells
form proper colorings of the experienced conflict relation.  It determines what
a successful representation-producing process must accomplish without
selecting a particular answer.

\section{Complete four-state enumeration}

The fifteen partitions of a four-element set have the following cell-size
patterns.
\begin{center}
\begin{tabular}{@{}lcc@{}}
\toprule
Cell-size pattern & Number & Role in the example\\
\midrule
$4$ & $1$ & Initial awareness\\
$3+1$ & $4$ & Singleton-exception refinements\\
$2+2$ & $3$ & Rows, columns, and diagonals\\
$2+1+1$ & $6$ & Three-cell refinements\\
$1+1+1+1$ & $1$ & Complete awareness\\
\midrule
Total & $15$ & \\
\bottomrule
\end{tabular}
\end{center}

The three balanced two-cell partitions are
\begin{align}
\Prow
&=
\big\{\{x_{00},x_{01}\},\{x_{10},x_{11}\}\big\},
\label{eq:row-partition}\\
\Pcol
&=
\big\{\{x_{00},x_{10}\},\{x_{01},x_{11}\}\big\},
\label{eq:column-partition}\\
\Pdiag
&=
\big\{\{x_{00},x_{11}\},\{x_{01},x_{10}\}\big\}.
\label{eq:diagonal-partition}
\end{align}
For $j,k\in\{0,1\}$, define the four singleton-exception partitions
\begin{equation}
\mathcal P^{jk}
=
\big\{\{x_{jk}\},X\setminus\{x_{jk}\}\big\}.
\label{eq:singleton-partition}
\end{equation}
These seven two-cell partitions are exactly the immediate refinements of
$\mathcal P^0=\{X\}$.

\begin{proposition}[Four minimally restorative refinements]
\label{prop:four-restorations}
For the discovery dataset $\overline D$ in \cref{eq:discovery-data},
\begin{equation}
\mathcal R(\mathcal P^0,\overline D)
=
\{\Prow,\Pcol,\mathcal P^{00},\mathcal P^{11}\}.
\label{eq:restorative-set}
\end{equation}
The uniquely active theories under these partitions, written in the state order
$(x_{00},x_{01},x_{10},x_{11})$, are
\begin{align*}
m^r&=(0,0,1,1),\\
m^c&=(0,1,0,1),\\
m^{00}&=(0,1,1,1),\\
m^{11}&=(0,0,0,1).
\end{align*}
Under the objective law in \cref{eq:true-row-law}, $m^r=d$ and the other three
formulations are false outside the discovery data.
\end{proposition}

\begin{proof}
By \cref{prop:conflict-characterization}, a restorative partition must separate
$x_{00}$ from $x_{11}$.  Exactly four two-cell partitions do so: rows, columns,
the partition isolating $x_{00}$, and the partition isolating $x_{11}$.  The
diagonal partition and the singleton partitions isolating $x_{01}$ or $x_{10}$
leave $x_{00}$ and $x_{11}$ in the same cell.  Every restorative partition with
three or four cells strictly refines at least one of the four restorative
two-cell partitions and is therefore not minimal.

In each of the four listed partitions, one cell contains the observation
$x_{00}\mapsto0$ and the other contains $x_{11}\mapsto1$.  Cell constancy
therefore determines the displayed theory uniquely.  Direct comparison with
$d=(0,0,1,1)$ establishes the final claim.
\end{proof}

\begin{proposition}[Multiplicity with one conflict]
\label{prop:multiplicity}
Let $C$ be the only non-conflict-free cell of the current partition, let
$|C|=q\geq2$, and suppose exactly one conflict edge has both endpoints in $C$.
There are $2^{q-2}$ minimally restorative refinements that replace $C$ by two
cells and leave every other awareness cell unchanged.
\end{proposition}

\begin{proof}
Fix the endpoint labeled zero in the first new cell and the endpoint labeled one
in the second.  Each of the remaining $q-2$ states can independently be placed
in either cell, producing $2^{q-2}$ binary partitions.  Fixing the two labeled
endpoints removes the duplication generated by exchanging cell names.  Every
such partition separates the sole conflict edge and is an immediate refinement
obtained by splitting $C$; hence it is minimally restorative.
\end{proof}

For $q=4$, \cref{prop:multiplicity} gives the four partitions in
\cref{prop:four-restorations}.  The result shows that restoration and
minimality direct inquiry but generally do not select its output.

\section{A structured but fallible representation-producing correspondence}

The partition lattice treats every binary refinement as equally immediate.  To
distinguish a one-coordinate factor from a two-coordinate interaction or
exception, the representation-producing process requires some comparison or
descriptive structure.

Define the analyst-level coordinate functions
\[
\rho(x_{rc})=r,
\qquad
\chi(x_{rc})=c.
\]
These functions describe analyst-level capacities of the non-deliberative
generative system to compare retained particulars.  They are not initially
elements of the firms' subjective event algebra or information possessed for
deliberation.  Before a representation is produced, the firm cannot name a row
or column predicate or optimize over the corresponding partitions.

For any binary partition $\mathcal P$, let $b_{\mathcal P}:X\to\{0,1\}$ be an
indicator of either one of its cells.  Coordinate $\rho$ is essential to
$b_{\mathcal P}$ if changing the row while holding the column fixed changes the
indicator for some column.  Define essentiality of $\chi$ analogously.
Complementing $b_{\mathcal P}$ does not alter its essential coordinates.

\begin{definition}[Coordinate complexity]
The coordinate complexity $\kappa(\mathcal P)$ of a binary partition is the
number of coordinates in $\{\rho,\chi\}$ essential to an indicator generating
that partition.
\end{definition}

It follows that
\begin{equation}
\kappa(\Prow)=\kappa(\Pcol)=1,
\label{eq:factor-complexity}
\end{equation}
whereas
\begin{equation}
\kappa(\Pdiag)=\kappa(\mathcal P^{jk})=2
\quad\text{for every }j,k\in\{0,1\}.
\label{eq:interaction-complexity}
\end{equation}
Rows and columns are one-factor distinctions.  The diagonal partition is an
interaction, while a singleton exception requires a conjunction of the two
coordinates.

\begin{definition}[Complexity-filtered restorative correspondence]
For any $(\mathcal P,D)$ whose minimally restorative refinements are binary,
define
\begin{equation}
\mathcal I(\mathcal P,D)
=
\underset{\mathcal P'\in\mathcal R(\mathcal P,D)}{\arg\min}
\;\kappa(\mathcal P').
\label{eq:inquiry-correspondence}
\end{equation}
\end{definition}

\begin{corollary}[Factor-based candidates]
\label{cor:factor-candidates}
For the discovery data,
\[
\mathcal I(\mathcal P^0,\overline D)=\{\Prow,\Pcol\}.
\]
\end{corollary}

\begin{proof}
Apply \cref{eq:factor-complexity,eq:interaction-complexity} to the four
partitions in \cref{eq:restorative-set}.
\end{proof}

The correspondence is not the firm's subjective answer menu and $\kappa$ is
not a cost the unaware firm consciously minimizes.  It is the analyst's model
of which conceptual products can be generated from the experienced contrast by
a particular representational capacity.  The imposed grammar is substantive
and contestable: it favors simple factor explanations, even in a world where a
conjunctive exception is objectively correct.  Alternative inquiry
technologies can be studied by varying the available grammar and its complexity
ordering.

\begin{definition}[Representation-producing technology]
\label{def:representation-technology}
A deterministic representation-producing technology for firm $i$ is an
analyst-level function
\[
\phi_i:\Pi(X)\times H\longrightarrow\Pi(X).
\]
The output $\phi_i(\mathcal P,h)=\mathcal P$ means that no new representation
is produced.  A strict output $\mathcal P\prec\phi_i(\mathcal P,h)$ is a
successful production event.  The technology is \emph{restorative} if every
strict output lies in $\mathcal R(\mathcal P,D(h))$.  When
$\mathcal I(\mathcal P,D(h))$ is defined, it is \emph{factor-seeking} if every
strict output lies in $\mathcal I(\mathcal P,D(h))$.
\end{definition}

The technology is not controlled through subjective optimization.  Its input
$h$ is the analyst's representation of the retained particulars on which the
non-deliberative generative system can operate.  Its output, and even the set of
outputs it could have produced, is not a pre-insight subjective menu.

\begin{assumption}[Dissonance-triggered production]
\label{ass:dissonance-trigger}
If $\widehat M(\mathcal P,D(h))\neq\emptysetset$, then
$\phi_i(\mathcal P,h)=\mathcal P$.  If the active set is empty, the technology
may fail, again producing $\mathcal P$, or successfully produce a minimally
restorative refinement.  The main factor-seeking specification further
restricts a successful output to $\mathcal I(\mathcal P,D(h))$.  If production
fails while the active set is empty, the firm carries forward its previously
committed policy until a successful production event occurs.  All results that
condition on an acquired refinement condition on such success.
\end{assumption}

The function $\phi_i$ has no direct objective-law argument.  Two causal worlds
that generate the same retained history and current awareness therefore
produce the same output.  Agent-specific prior understanding, analogy, and
organizational routines may be represented by differences across the functions
$\phi_i$, but not by allowing a function to inspect unobserved values of the
objective law.

\section{Nonidentification and the cost of robustness}

The focal truth is the row law $d^r\equiv d$, where $d^r(x_{rc})=r$.  For comparison, define the
column law
\begin{equation}
d^c(x_{rc})=c.
\label{eq:column-law}
\end{equation}
Both causal worlds agree on the discovery states:
\[
d^r(x_{00})=d^c(x_{00})=0,
\qquad
d^r(x_{11})=d^c(x_{11})=1.
\]

\begin{definition}[Truth-blind inquiry-and-formulation rule]
A single-valued inquiry-and-formulation rule returns, conditional on successful
representation production, one pair $(\mathcal P',m')$ satisfying
\[
\mathcal P'=\phi_i(\mathcal P,h),
\qquad
m'\in\widehat M(\mathcal P',D(h)).
\]
It is truth-blind if two objective causal laws that generate the same current
awareness partition and retained history induce the same pair.  The rule is
correct under objective law $d'$ if $m'=d'$.  It may depend on experience,
retained particulars, the comparison grammar, and agent-specific background
capacities, but not on unobserved values of the objective law.
\end{definition}

\begin{theorem}[Observational nonidentification]
\label{thm:nonidentification}
No deterministic truth-blind inquiry-and-formulation rule can produce the correct causal
formulation after $\overline D$ in both the row world $d^r$ and the column world
$d^c$.
\end{theorem}

\begin{proof}
The current partition and retained history are identical in the two worlds, so
a deterministic truth-blind rule must produce the same formulation in both.
But $d^r\neq d^c$.  The same formulated theory cannot equal both objective
laws.
\end{proof}

This result does not depend on the coordinate-complexity criterion.  It follows
from observational equivalence.  The $2\times2$ geometry provides an additional
symmetry interpretation.  Let
\[
\pi(x_{rc})=x_{cr}
\]
exchange the coordinates.  The transformation fixes $x_{00}$, $x_{11}$,
$\mathcal P^0$, and $\overline D$, while exchanging $\Prow$ and $\Pcol$.
For a partition selector $s$, coordinate equivariance means
\[
s(\pi\mathcal P,\pi h)=\pi s(\mathcal P,h).
\]
Here $\pi\mathcal P$ and $\pi h$ apply the state relabeling to every partition
cell and every customer-state entry in the retained history.
Therefore no deterministic coordinate-equivariant selector can choose one of
the two elements of $\mathcal I(\mathcal P^0,\overline D)$ without an
additional asymmetry in evidence or background accessibility.  The complexity
filter, rather than symmetry alone, is what excludes the invariant singleton-
exception candidates.

\begin{definition}[Causal sufficiency]
Partition $\mathcal P$ is sufficient for a causal law $d'$ if
\[
d'\in M(\mathcal P),
\]
or, equivalently, if $d'$ is constant on every cell of $\mathcal P$.
\end{definition}

\begin{corollary}[Worst-case success of minimal inquiry]
\label{cor:maximin}
For comparison with the deterministic baseline, suppose an analyst-level
truth-blind technology randomizes over
$\mathcal R(\mathcal P^0,\overline D)$, and call it successful in a causal world
when its selected partition is sufficient to express that world's objective
law.  Randomized truth blindness requires observationally equivalent causal
worlds to induce the same distribution over outputs.  Its worst-case
probability of success across $\{d^r,d^c\}$ is at most
$1/2$.  The bound is attained by assigning equal probability to $\Prow$ and
$\Pcol$.
\end{corollary}

\begin{proof}
Among the four minimally restorative partitions, only $\Prow$ is sufficient
for $d^r$ and only $\Pcol$ is sufficient for $d^c$.  If their selection
probabilities are $p_r$ and $p_c$, then $p_r+p_c\leq1$, so
$\min\{p_r,p_c\}\leq1/2$.
\end{proof}

\begin{proposition}[Robust sufficiency requires complete awareness]
\label{prop:robust-awareness}
If a partition is sufficient for both $d^r$ and $d^c$, then it is the singleton
partition
\begin{equation}
\Pfull
=
\big\{
\{x_{00}\},\{x_{01}\},\{x_{10}\},\{x_{11}\}
\big\}.
\label{eq:full-partition}
\end{equation}
\end{proposition}

\begin{proof}
Sufficiency for $d^r$ requires every cell to lie within a row.  Sufficiency for
$d^c$ requires every cell to lie within a column.  Each row-column intersection
contains exactly one state.  Hence every cell is a singleton.
\end{proof}

\Cref{prop:robust-awareness} concerns representation, not identification.
Complete awareness makes both observationally equivalent causal laws
expressible, but $\overline D$ alone still does not determine which one is true.
Thus exhaustive awareness is necessary for a representation guaranteed to
accommodate both worlds, but is not sufficient for knowledge.  A non-exhaustive
technology avoids brute-force bookkeeping only by accepting fallibility.

\section{Formulation, representation-relative coherence, and judgment}

Successful representation production first supplies a new partition
$\mathcal P'$.  A causal formulation is then a pair
\[
(\mathcal P',m'),
\qquad
m'\in\widehat M(\mathcal P',\overline D).
\]
In the focal example, every minimally restorative partition in
\cref{eq:restorative-set} has a unique active theory, so representation and
formulation occur in immediate succession.  They remain conceptually distinct:
the partition supplies a distinction; the theory states the causal relation
organized by that distinction.

After acquiring $\mathcal P'$, firm $i$ forms a bridge prior
\begin{equation}
\mu_i^+\in\Delta(M(\mathcal P')).
\label{eq:bridge-prior}
\end{equation}
Ordinary Bayesian updating does not derive $\mu_i^+$ from $\mu_i^0$, because
the domains differ.  For the deterministic example, assume only that the bridge
prior has full support.  Define the representation-relative coherence
distribution
\begin{equation}
\overline\mu_i(m\mid\mathcal P',\overline D)
=
\frac{
\mu_i^+(m)\ind\{m\in\widehat M(\mathcal P',\overline D)\}
}{
\sum_{m'\in M(\mathcal P')}
\mu_i^+(m')\ind\{m'\in\widehat M(\mathcal P',\overline D)\}
}.
\label{eq:coherence-distribution}
\end{equation}
This is a retrospective renormalization on the newly produced model space, not
a posterior generated from one temporally prior probability model over the
entire representation-selection process.

\begin{proposition}[Representation-relative certainty under a false formulation]
\label{prop:false-coherence}
Suppose inquiry produces $\Pcol$ and $\mu_i^+$ has full support on
$M(\Pcol)$.  The coherence distribution assigns probability one to
\[
m^c=(0,1,0,1).
\]
Nevertheless, under the row truth,
\[
m^c(x_{01})\neq d^r(x_{01}),
\qquad
m^c(x_{10})\neq d^r(x_{10}).
\]
\end{proposition}

\begin{proof}
Under $\Pcol$, cell constancy and $x_{00}\mapsto0$ require value zero on the
first column.  Cell constancy and $x_{11}\mapsto1$ require value one on the
second.  Hence $m^c$ is the unique model in
$\widehat M(\Pcol,\overline D)$, and full support gives it
$\overline\mu_i$-probability one.  Direct evaluation at the off-diagonal states
establishes both errors.
\end{proof}

The distribution in \cref{prop:false-coherence} is Bayesian-coherent
conditional on the generated representation.  It should not be treated as
independent confirmation: the same evidence directed representation production
and fitted the resulting theory.  To implement reflective judgment cleanly in
this model, evidential assessment must use fresh evidence or a likelihood that
explicitly accounts for the representation-selection process.  The remainder
of this section makes that accounting exact.

\begin{definition}[Reflective assessment problem]
\label{def:reflective-assessment}
A reflective assessment problem is a pair $(\mathcal D,\nu)$, where
$\mathcal D$ is a finite set of candidate objective laws $d':X\to A$ and
$\nu\in\Delta(\mathcal D)$ is a full-support reflective prior.  The assessment
asks how much probability the total evidence---the retained history together
with the record of what inquiry produced---places on each candidate law.
\end{definition}

\begin{definition}[Selection-adjusted warrant and the warrant gap]
\label{def:warrant}
Fix a history $h$ and a successful production event yielding the formulation
$(\mathcal P',m')$.  The \emph{selection-adjusted warrant} of the formulation
is
\begin{equation}
W(m'\mid h)
=
\nu\big(\{d\in\mathcal D:d=m'\}\;\big|\;
h,\ \phi_i(\mathcal P,h)=\mathcal P'\big),
\label{eq:warrant}
\end{equation}
the posterior probability that the objective law coincides with the formulated
theory, conditional on both the retained history and the production event.
The \emph{warrant gap} is
\begin{equation}
G(m'\mid h)
=
\overline\mu_i(m'\mid\mathcal P',D(h))-W(m'\mid h).
\label{eq:warrant-gap}
\end{equation}
\end{definition}

\begin{lemma}[Guaranteed fit]
\label{lem:guaranteed-fit}
Conditional on a successful production event, the produced partition lies in
$\mathcal R(\mathcal P,D(h))$, so its active set is nonempty and every active
theory agrees with all revealed labels.  The event that the formulation fits
the discovery data therefore has probability one under every candidate law
that generates $h$, and retrospective fit carries no discriminating
information.
\end{lemma}

\begin{proof}
Immediate from \cref{def:representation-technology} and \cref{eq:active-set}:
a restorative output is defined by the nonemptiness of its active set, and
membership of the active set is defined by agreement with $D(h)$.  Neither
condition involves the objective law except through $D(h)$ itself.
\end{proof}

\begin{proposition}[Prior preservation under truth-blind production]
\label{prop:prior-preservation}
Suppose the firms' committed policies are fixed, the technology is truth-blind
(deterministic or randomized), and each candidate law consistent with the
revealed labels generates the same realized history $h$.  Then conditioning on
$h$ and on the production event excludes the candidates inconsistent with the
revealed labels and preserves the reflective prior odds among the rest.  For
the discovery history with $\mathcal D=\{d^r,d^c\}$ and symmetric $\nu$, both
candidates survive; production of either factor-based refinement yields a
coherence probability of one on its unique active theory
(\cref{prop:false-coherence}); and
\[
W(m^r\mid h)=W(m^c\mid h)=\tfrac12,
\qquad
G(m^r\mid h)=G(m^c\mid h)=\tfrac12.
\]
\end{proposition}

\begin{proof}
The reflective posterior on $d\in\mathcal D$ is proportional to the product of
three factors: the prior $\nu(d)$; the likelihood of the realized history,
which in the deterministic environment is an indicator equal to one exactly
when $d$ agrees with every label in $D(h)$, since a disagreeing law would have
produced different payoffs at some differentiated episode; and the probability
of the production event given $(d,h)$, which by truth blindness depends on
$(\mathcal P,h)$ alone and is therefore constant across $\mathcal D$ and
cancels.  This establishes the general claim.  For the discovery history,
$d^r$ and $d^c$ agree with both labels in $\overline D$
(\cref{thm:nonidentification}), so the symmetric prior is preserved:
$\nu(d^c\mid\cdot)=\nu(d^r\mid\cdot)=1/2$.  Since $m^c=d^c\neq d^r$, the
warrant of $m^c$ is one-half, and symmetrically for $m^r$.  Subtracting from
the unit coherence probability gives the gap.
\end{proof}

\begin{corollary}[Self-audit requires complete awareness]
\label{cor:self-audit}
The reflective assessment $(\mathcal D,\nu)$ is expressible by firm $i$ only
if every candidate law in $\mathcal D$ is a theory expressible under
$\mathcal P_i$, that is, only if $\mathcal P_i$ is sufficient for every
$d\in\mathcal D$.  For $\mathcal D=\{d^r,d^c\}$,
\cref{prop:robust-awareness} implies $\mathcal P_i=\Pfull$.  A firm holding
either minimally restorative factor candidate therefore cannot formulate the
comparison that measures its own warrant gap: the overstatement identified in
\cref{prop:false-coherence} is invisible from within the representation that
produces it.  Short of complete awareness, warrant must be supplied
externally---by fresh diagnostic evidence, an analyst, a differently aware
rival, or a future self with expanded awareness.
\end{corollary}

\begin{proof}
Holding $\nu$ over $\mathcal D$ requires entertaining each candidate law as a
hypothesis about customer response, which requires each law to be cell
constant under $\mathcal P_i$ (\cref{sec:unawareness}); this is causal
sufficiency for each $d\in\mathcal D$.  Sufficiency for both $d^r$ and $d^c$
forces the singleton partition by \cref{prop:robust-awareness}.
\end{proof}

\begin{remark}[Double use, made exact]
\label{rem:double-use}
The coherence distribution treats restored fit as confirmation;
\cref{lem:guaranteed-fit} shows the fit was guaranteed by the restoration
property, and \cref{prop:prior-preservation} shows the production event
carries no additional information when the technology is truth-blind.  The
warrant gap of one-half is therefore an exact measure of the confidence
overstatement produced by using the discovery experience twice---once to
direct representation production and once to fit the resulting theory.  Truth
blindness is the polar case in which the production event is completely
uninformative.  If production were correlated with the objective law---for
instance, social transmission from an informed source---the event would itself
be evidence, and the selection adjustment could raise as well as lower
confidence.
\end{remark}

\begin{remark}[Unwarranted certainty is not a symptom of error]
\label{rem:warranted-luck}
The gap afflicts the correct formulation equally: the firm that acquires
$\Prow$ holds $\overline\mu=1$ on $m^r=d^r$, yet its warrant at the discovery
stage is also one-half.  Early certainty is unwarranted even when it happens
to be true.  What separates the two firms is not their evidence or their
coherence but the objective law, which neither observes; validation by fresh
differentiated evidence is what converts a fortunate formulation into
warranted knowledge.
\end{remark}

The row and column formulations make opposing off-diagonal predictions:
\begin{center}
\begin{tabular}{@{}lcccc@{}}
\toprule
Formulation & $x_{00}$ & $x_{01}$ & $x_{10}$ & $x_{11}$\\
\midrule
Row $m^r$ & $0$ & $0$ & $1$ & $1$\\
Column $m^c$ & $0$ & $1$ & $0$ & $1$\\
\bottomrule
\end{tabular}
\end{center}
Either off-diagonal state can therefore discriminate between them, provided
strategic actions reveal its objective label.

\begin{proposition}[Diagnostic closure of the warrant gap]
\label{prop:warrant-closure}
Append additional episodes to the discovery history and maintain the
reflective assessment of \cref{prop:prior-preservation}.
\begin{enumerate}[label=(\roman*)]
  \item A single differentiated episode at an off-diagonal state closes the
  gap.  If the revealed label agrees with the acquired formulation, $W$ rises
  to one and $G=0$.  If it disagrees, the active set under the acquired
  partition empties, coherence collapses with renewed dissonance, and the
  warrant of the surviving candidate rises to one.
  \item Pooled episodes, at any states and in any number, change neither the
  coherence distribution nor the selection-adjusted posterior: the warrant gap
  is invariant under unbounded pooled experience.
\end{enumerate}
\end{proposition}

\begin{proof}
(i) At $x_{01}$ the candidates predict opposite labels ($d^r$ assigns $0$ and
$d^c$ assigns $1$), and likewise at $x_{10}$.  A revealed label therefore has
likelihood one under one candidate and zero under the other, making the
reflective posterior degenerate.  If the label agrees with the acquired
formulation, the active set is unchanged, $\overline\mu$ remains one, and
$G=1-1=0$.  If it disagrees, the new label conflicts with a previously
revealed diagonal label inside one cell of the acquired partition, the active
set empties, and \cref{ass:dissonance-trigger} governs renewed inquiry.
(ii) By \cref{lem:revelation}, a pooled episode has the same payoff vector
under every candidate law, so its likelihood is one under both candidates and
the reflective posterior is unchanged; the active set is also unchanged, so
$\overline\mu$ is unchanged.
\end{proof}

\begin{proposition}[Resolution under one additional diagnostic label]
\label{prop:directed-resolution}
Suppose representation production initially fails after $\overline D$, so
awareness remains $\mathcal P^0$ and the previously differentiated policy
profile is carried forward for one additional diagnostic episode.  Let the row
law be true.
\begin{enumerate}[label=(\alph*)]
  \item If differentiated actions reveal $x_{01}\mapsto0$, the minimally
  restorative refinements are $\Prow$ and $\mathcal P^{11}$, and the unique
  minimum-complexity refinement is $\Prow$.
  \item If differentiated actions reveal $x_{10}\mapsto1$, the minimally
  restorative refinements are $\Prow$ and $\mathcal P^{00}$, and the unique
  minimum-complexity refinement is $\Prow$.
\end{enumerate}
\end{proposition}

\begin{proof}
For part (a), the two zero-labeled states must lie apart from $x_{11}$.  In a
minimal binary restoration, the unobserved state $x_{10}$ may join either side,
giving the row partition and the partition isolating $x_{11}$.  Their
complexities are one and two.  Part (b) is symmetric.
\end{proof}

The proposition is a continued-commitment comparative exercise rather than a
claim that undefined post-dissonance beliefs generate an equilibrium policy.
It is conditional on obtaining a fresh diagnostic observation.  An
unaware firm cannot simply be assumed to target ``the off-diagonal state'' as
an experiment before it can formulate that distinction.  A theory of inquiry
must explain whether the observation arrives through ordinary experience,
comparison of retained cases, social transmission, or an exploration
technology available without prior semantic specification.

\section{Cumulative awareness and path dependence}

\begin{assumption}[Cumulative awareness]
\label{ass:cumulative-awareness}
Once a distinction becomes available, later awareness can refine but does not
coarsen the acquired partition.  Rejecting a causal theory need not make the
distinction on which it was formulated unthinkable again.
\end{assumption}

This is an explicit economic assumption, not a logical implication of
unawareness.  It permits causal theory to be rejected while the associated
conceptual distinction remains available.

\begin{proposition}[Path-dependent terminal awareness]
\label{prop:path-dependence}
Assume the row law, cumulative awareness, successful minimally restorative
revision after every dissonance, and eventual differentiated observations at
both off-diagonal states.
\begin{enumerate}[label=(\roman*)]
  \item A firm that initially acquires $\Prow$ holds the correct formulation
  $m^r=d^r$ and need not refine further.
  \item A firm that initially acquires $\Pcol$ must eventually reach
  $\Pfull$.
\end{enumerate}
The two firms can therefore converge to the same correct causal theory and
policy while retaining different awareness partitions.
\end{proposition}

\begin{proof}
Under $\Prow$, every subsequently revealed label agrees with $m^r$, so the
active set never becomes empty.

Under $\Pcol$, $x_{00}$ and $x_{10}$ initially share one cell, while $x_{01}$
and $x_{11}$ share the other.  A revealed label at $x_{01}$ conflicts with the
previous label at $x_{11}$ and forces the second column to split.  A revealed
label at $x_{10}$ conflicts with the previous label at $x_{00}$ and forces the
first column to split.  Reversing the order of the off-diagonal observations
reverses the order of these two splits.  Cumulative awareness rules out
discarding the column distinction and replacing it with the incomparable row
partition.  The terminal partition is therefore the coarsest common refinement
\[
\Pcol\vee\Prow
=
\Pfull.
\]
Once all four labels have been revealed, the unique active causal function is
$d^r$ under either awareness path.
\end{proof}

\begin{remark}[Alternative false paths]
If the complexity filter is not imposed, an initial singleton formulation also
leaves a permanent representational footprint.  Under complete diagnostic
coverage, $\mathcal P^{00}$ leads to
\[
\big\{\{x_{00}\},\{x_{01}\},\{x_{10},x_{11}\}\big\},
\]
whereas $\mathcal P^{11}$ leads to
\[
\big\{\{x_{00},x_{01}\},\{x_{10}\},\{x_{11}\}\big\}.
\]
Both are sufficient for the row law but finer than the minimally sufficient
row partition.
\end{remark}

\begin{proposition}[Diagnostic-coverage benchmark]
\label{prop:coverage}
Let $X$ be finite and $d:X\to\{0,1\}$ deterministic.  Suppose awareness only
changes after dissonance through strict refinements; after every dissonance,
revision continues until a partition with a nonempty active set is produced;
and every $x\in X$ is eventually observed under differentiated actions.  Let
the terminal partition be the partition in place after the last revealing
observation and any revision it requires.  Then the terminal partition exists,
every terminal awareness cell is homogeneous under $d$, and the terminal active
model set is $\{d\}$.
\end{proposition}

\begin{proof}
Every successful awareness change strictly increases the number of cells, so
finiteness of $X$ permits only finitely many such changes.  Hence a terminal
partition exists after the final revealing observation and required revision.
Differentiated observation reveals $d(x)$ for every state.  An active theory
must consequently satisfy $m(x)=d(x)$ for all $x\in X$, so it equals $d$.
No terminal awareness cell can contain two states with different objective
labels, because cell constancy would then make the active set empty.
\end{proof}

\Cref{prop:coverage} is a benchmark rather than a headline result.  Its strong
coverage condition does most of the work.  The economically substantive issue
is whether endogenous strategy supplies or suppresses differentiated
observations.

\section{Strategic diagnosticity and self-confirming awareness}

The policies analyzed in this section are provisional commitments based on
representation-relative beliefs.  They may themselves generate validation
evidence; they are not assumed to follow a completed reflective judgment.

Once a partition is acquired prospectively, a firm with a degenerate belief on
$m$ follows the policy
\[
\sigma_i(C)=m(C)
\]
at every cell labeled zero or one.  By \cref{lem:dominance}, this policy is a
subjectively strict best response against either rival action.

\begin{definition}[Self-confirming awareness trap]
\label{def:awareness-trap}
A profile of awareness partitions, beliefs, and policies forms a
self-confirming awareness trap if
\begin{enumerate}[label=(\roman*)]
  \item every firm's policy is optimal under its belief;
  \item all outcomes generated by the policy profile are consistent with each
  firm's active causal theory, so no dissonance triggers further inquiry; and
  \item there exist a firm $i$ with active theory $m_i$ and a state $x$ such
  that $m_i(x)\neq d(x)$, but the on-path policy profile pools at $x$ and hence
  does not reveal $d(x)$.
\end{enumerate}
\end{definition}

\begin{theorem}[Shared false formulation suppresses its own falsification]
\label{thm:self-confirming-trap}
Suppose the row law is true, both firms acquire $\Pcol$, both become certain of
$m^c=(0,1,0,1)$ after the discovery batch, and every customer state subsequently
recurs.  Suppose \cref{ass:dissonance-trigger} governs later representation
production and the firms recommit their optimal policies in every generation.
The induced policy profile is a self-confirming awareness trap.
The false formulation is never falsified, even though it is wrong at both
off-diagonal states.
\end{theorem}

\begin{proof}
Under certainty of $m^c$, \cref{lem:dominance} makes the column policy
\[
\sigma_i(x_{00},x_{01},x_{10},x_{11})=(0,1,0,1)
\]
subjectively strictly optimal for each firm.  The firms therefore choose the
same action at every state.  By \cref{lem:revelation}, every subsequent payoff
vector is $(1/2,1/2)$ and supplies no causal label.  The pooled payoffs are
exactly those predicted by $m^c$ under the on-path action profile, so the active
set never becomes empty.  Nevertheless, $m^c(x_{01})=1\neq0=d^r(x_{01})$ and
$m^c(x_{10})=0\neq1=d^r(x_{10})$.
\end{proof}

The trap relies on the market-splitting payoff in \cref{eq:objective-payoff}:
pooled firms receive one-half even when both offer the objectively inferior
architecture.  Introducing an outside option or demand losses could make an
incorrect pooled architecture directly informative and would alter the result.

\begin{corollary}[Permanent warrant gap at the trap]
\label{cor:permanent-gap}
In the environment of \cref{thm:self-confirming-trap}, every post-discovery
episode is pooled, so by \cref{prop:warrant-closure}(ii) the warrant gap
remains $G=\tfrac12$ in every generation.  The firms accumulate unbounded
experience, remain representation-relatively certain of $m^c$, and are never
more than half-warranted; the missing half is exactly the reflective weight on
the causal world their shared partition cannot express.
\end{corollary}

\begin{proof}
Immediate from \cref{thm:self-confirming-trap} and
\cref{prop:warrant-closure}(ii).
\end{proof}

\begin{corollary}[Theory heterogeneity creates diagnostic competition]
\label{cor:heterogeneous-theories}
Suppose the row law is true, firm $1$ is certain of $m^r$, and firm $2$ is
certain of $m^c$.  Their policies differ at $x_{01}$ and $x_{10}$, revealing
the objective labels and falsifying $m^c$.  If all four states are served once
with equal weight in one committed validation batch before either firm can
revise, firm $1$ earns average payoff $3/4$ and firm $2$ earns $1/4$.
\end{corollary}

\begin{proof}
The policies are
\[
\sigma_1=(0,0,1,1),
\qquad
\sigma_2=(0,1,0,1).
\]
They pool at $x_{00}$ and $x_{11}$, giving each firm one-half.  At $x_{01}$ and
$x_{10}$, firm $1$ offers the objectively preferred architecture and receives
one.  Thus its equally weighted payoff is
\[
\frac14\left(\frac12+1+1+\frac12\right)=\frac34,
\]
with the complementary payoff accruing to firm $2$.  Both differentiated
outcomes reveal the labels on which $m^c$ is wrong.
\end{proof}

The theorem and corollary identify an endogenous epistemic effect of strategic
interaction.  Shared theories align actions and can suppress causal evidence.
Different theories create action differentiation, which both reallocates value
capture and generates the evidence needed for judgment.

\section{The economics of inquiry: costly activation and delabeled leverage}
\label{sec:inquiry-economics}

In the preceding sections, representation production is triggered by
dissonance and costs nothing, as if activation were free and automatic.  A
theory of directed inquiry requires more: whether to prosecute a question must
itself be a decision.  The obstacle is that the natural criterion for that
decision---the expected value of the answer---is doubly unavailable at exactly
the moment the question exists.  The Bayesian posterior is undefined
(\cref{lem:bayesian-closure}), and the candidate answers are not objects of
pre-insight deliberation (\cref{def:representation-technology}).  This section
shows that the laboratory nevertheless supports an honest valuation: a
statistic computed entirely from retained episodes and revealed payoffs that
is invariant across every representation successful inquiry could produce.

\begin{assumption}[Costly activation]
\label{ass:costly-activation}
At the close of a generation whose history exhibits causal dissonance, firm
$i$ elects an inquiry effort $\varepsilon_i\in\{0,1\}$ at cost
$\gamma_i\varepsilon_i$, with $\gamma_i>0$ in payoff units.  The technology
$\phi_i$ of \cref{def:representation-technology} operates only if
$\varepsilon_i=1$ and then succeeds or fails as in
\cref{ass:dissonance-trigger}.  If $\varepsilon_i=0$ or production fails, the
firm carries forward its previously committed policy.  Firm $i$ holds a
delabeled success weight $\theta_i\in(0,1]$: a subjective probability that an
elected inquiry episode produces a usable refinement.  The weight is defined
on the event of success alone; it is not a distribution over the identities of
possible outputs.
\end{assumption}

The success weight is awareness of unawareness in delabeled form.  The firm
can believe that inquiry may yield a usable distinction without being able to
describe any candidate distinction.  The earlier sections are the special case
$\gamma_i=0$, in which election is always subjectively justified and
activation is effectively automatic.

\begin{definition}[Revealed labels and the disagreement set]
\label{def:disagreement-set}
A retained episode $s$ in $e_i(h;\mathcal P_i)$ is \emph{labeled} if
$a_{1s}\neq a_{2s}$; its label $a^*_s$ is the architecture of the firm that
received payoff one.  Let $\sigma_i$ denote firm $i$'s currently committed
policy.  The \emph{disagreement set} is
\[
L_i(h)
=
\left\{
s:\ s\text{ labeled and }
\sigma_i\big(C_{\mathcal P_i}(x_s)\big)\neq a^*_s
\right\},
\]
the labeled episodes at whose recurrence the committed policy prescribes the
architecture that lost.  The \emph{retained-record leverage} is
\begin{equation}
B_i(h)=\tfrac12\,\bigl|L_i(h)\bigr|.
\label{eq:leverage}
\end{equation}
\end{definition}

Both objects are deliberative-level computations.  A token's awareness cell,
the actions taken, and the payoffs received are all recorded in
$e_i(h;\mathcal P_i)$, so identifying the winning architecture requires no
predicate the firm lacks.  For a firm with coarse awareness and a constant
committed policy, $L_i(h)$ is exactly the set of differentiated episodes the
firm lost.

\begin{assumption}[Subjective recurrence and pinned valuation]
\label{ass:subjective-recurrence}
In evaluating inquiry, firm $i$ proceeds as if the next generation's batch
reproduces its retained record token for token, and it evaluates one
generation ahead.  It values a prospective refinement by its \emph{pinned
gain}: the component of the anticipated payoff change that is common to every
restorative refinement.  The recurrence premise is a belief, not foresight;
the analyst does not require the future batch to satisfy it, and a
miscalibrated premise can produce both under- and over-investment in inquiry.
A longer subjective horizon scales the valuation by the number of anticipated
recurrences and changes nothing qualitative.
\end{assumption}

\begin{lemma}[Unit leverage of a disagreeing episode]
\label{lem:unit-leverage}
Consider the recurrence of a retained labeled episode $s$ and compare two
policies for firm $i$: the carried-forward committed policy, and a policy
identical to it except that it implements the label $a^*_s$ at that
recurrence.  Against every rival action, pure or mixed:
\begin{enumerate}[label=(\roman*)]
  \item if $s\in L_i(h)$, the implementing policy gains exactly one-half;
  \item if $s\notin L_i(h)$, the two policies coincide at the recurrence and
  the gain is zero.
\end{enumerate}
\end{lemma}

\begin{proof}
For (i), write $a$ for the committed action at the recurrence and
$a^*=a^*_s\neq a$ for the label; by \cref{lem:revelation}, $a^*=d(x_s)$.  If
the rival plays $a^*$, the committed policy is differentiated and loses,
yielding zero, while the implementing policy pools, yielding one-half.  If the
rival plays $a$, the committed policy pools, yielding one-half, while the
implementing policy is differentiated and wins, yielding one.  In both cases
the difference is one-half, and linearity extends the conclusion to mixtures.
For (ii), a labeled episode outside $L_i(h)$ satisfies
$\sigma_i(C_{\mathcal P_i}(x_s))=a^*_s$, so implementation changes nothing.
\end{proof}

The gain in part (i) is independent of the rival's behavior, so no conjecture
about the rival is needed to price the disagreement set.  The structural
reason is the same one driving \cref{lem:dominance}: with two architectures
and payoffs summing to one, switching from the losing to the winning
architecture at a state moves the firm up by exactly one-half against either
rival action.

\begin{lemma}[Answer invariance]
\label{lem:answer-invariance}
Let the current partition exhibit causal dissonance at $D(h)$, let
$\mathcal P'\in\mathcal R(\mathcal P_i,D(h))$ be any minimally restorative
refinement, and let $m'\in\widehat M(\mathcal P',D(h))$ be held with
certainty.  Then the induced subjectively optimal policy implements $a^*_s$ at
the recurrence of every labeled episode $s$, and the anticipated pinned gain
over the carried-forward policy equals $B_i(h)$, identically across all
elements of $\mathcal R(\mathcal P_i,D(h))$ and their active theories.  At
unlabeled recurrences the induced policy may also depart from the
carried-forward policy; the value of that departure depends on unrevealed
values of the objective law and is not represented in $B_i(h)$.
\end{lemma}

\begin{proof}
By \cref{prop:conflict-characterization}, every cell of $\mathcal P'$ is
conflict-free, so by \cref{eq:active-set} every active theory satisfies
$m'(x_s)=a^*_s$ at each labeled episode.  Under certainty of $m'$,
\cref{lem:dominance} makes the policy $\sigma_i(C)=m'(C)$ subjectively
strictly optimal at every cell labeled zero or one, and
\cref{ass:recodability} licenses its prospective implementation.  Summing
\cref{lem:unit-leverage} over labeled recurrences gives a total pinned gain of
$\lvert L_i(h)\rvert/2=B_i(h)$, which does not depend on which restorative
refinement was acquired.
\end{proof}

\begin{remark}[Delabeled deliberation]
\label{rem:delabeled-deliberation}
The leverage computation quantifies over retained tokens and their recorded
outcomes, never over candidate representations.  The firm's reasoning can be
glossed: ``customers like the episodes I lost will recur; whatever distinction
inquiry may yield, possessing it would let me offer such customers the
architecture that won.''  The gloss uses only token-indexicals, which
retention supplies (\cref{ass:recodability}); it requires no predicate, no
answer menu, and no probability over semantic alternatives.  This is the
model's precise sense in which a question can be priced before its answer can
be thought.
\end{remark}

\begin{remark}[Pinned versus extrapolated value]
\label{rem:pinned-extrapolated}
The statistic $B_i$ prices only what every answer must deliver: agreement with
the revealed labels.  A full formulation extrapolates beyond them---for
example, the column formulation's off-diagonal prescriptions---and the
extrapolated component is exactly what representation-relative coherence
cannot warrant (\cref{prop:false-coherence}) and what reflective judgment must
assess with fresh or selection-adjusted evidence.  Anticipated leverage and
warranted confidence therefore separate cleanly: the former is
answer-invariant and computable at dissonance; the latter is answer-specific
and computable only after formulation.
\end{remark}

\begin{proposition}[Participation rule]
\label{prop:participation}
Under \cref{ass:costly-activation,ass:subjective-recurrence}, electing inquiry
has represented value $\theta_iB_i(h)-\gamma_i$ and abstention has represented
value zero, so firm $i$ elects inquiry if and only if
\[
\theta_iB_i(h)\geq\gamma_i,
\]
with indifference resolved toward inquiry.  After the discovery batch
$\overline D$, $\lvert L_i(h)\rvert=1$ for both firms, so each elects inquiry
if and only if $\theta_i/2\geq\gamma_i$.
\end{proposition}

\begin{proof}
Election costs $\gamma_i$ and, with subjective probability $\theta_i$,
production succeeds; by \cref{lem:answer-invariance} the pinned gain on the
recurrence record is then $B_i(h)$ whichever refinement arrives, and
\cref{ass:subjective-recurrence} assigns the extrapolated component no
represented value at the election.  On failure the carried-forward policy
continues and the gain is zero.  For the discovery batch, firm $1$'s constant
policy $0$ agrees with the label at $x_{00}$ and disagrees at $x_{11}$, while
firm $2$'s constant policy $1$ disagrees only at $x_{00}$; hence
$\lvert L_i\rvert=1$ and $B_i=1/2$ for both firms.
\end{proof}

Both firms hold identical leverage despite opposite experiences: each won one
episode and lost the other.  The incentive to inquire is symmetric; which firm
succeeds, and on which refinement, is governed by the technologies $\phi_i$,
the locus of firm-specific insight capacity.

\begin{proposition}[Action separation, not residual uncertainty]
\label{prop:action-separation}
$B_i(h)>0$ if and only if some retained labeled episode contradicts the
committed policy.  In particular:
\begin{enumerate}[label=(\roman*)]
  \item on the self-confirming path of \cref{thm:self-confirming-trap}, no
  episode is ever labeled, so $L_i(h)=\emptysetset$ and $B_i(h)=0$ for both
  firms in every generation, although each firm's formulation is false at two
  of the four states and the empty label record leaves the entire model space
  unrestricted at the analyst's level;
  \item after the discovery batch, the revealed labels reduce the minimally
  restorative refinements to four (\cref{prop:four-restorations}) and the
  complexity-filtered candidates to two
  (\cref{cor:factor-candidates}), yet $B_i(h)=1/2>0$.
\end{enumerate}
The represented value of a question therefore tracks the action separation in
the retained record, not the quantity of unresolved structure.
\end{proposition}

\begin{proof}
The equivalence restates \cref{def:disagreement-set}.  For (i), pooled actions
label no episode by \cref{lem:revelation}.  For (ii), the count is computed in
\cref{prop:participation}, and the reduction of candidates is given by the
cited results.
\end{proof}

\begin{corollary}[Trap immunity to cheap inquiry]
\label{cor:trap-immunity}
Extend the environment of \cref{thm:self-confirming-trap} with
\cref{ass:costly-activation}.  The trap persists for every $\gamma_i>0$ and
every $\theta_i\leq1$.  It is doubly locked: on-path experience generates no
dissonance, so the technology is inert (\cref{ass:dissonance-trigger}); and
$B_i(h)=0$ in every generation, so even an infallible technology at an
arbitrarily small positive cost would receive no represented justification.
Escaping the trap therefore requires a source of action differentiation
outside the represented calculus---exogenous experimentation, entry by a firm
with a different formulation, or a payoff environment in which pooled
misservice is itself informative.
\end{corollary}

\begin{proof}
Combine \cref{thm:self-confirming-trap} with
\cref{prop:action-separation}(i) and the participation rule of
\cref{prop:participation}.
\end{proof}

\begin{corollary}[Rational cessation and the loser's leverage]
\label{cor:loser-leverage}
Append the validation batch of \cref{cor:heterogeneous-theories} to the
discovery history, with firm $1$ certain of $m^r$ and firm $2$ certain of
$m^c$.
\begin{enumerate}[label=(\roman*)]
  \item Firm $1$'s record contains labels at all four states, each agreeing
  with its committed policy: $L_1=\emptysetset$ and $B_1=0$.  Firm $1$
  experiences no dissonance and, even were inquiry offered, would decline it
  at every positive cost.  Its terminal awareness $\Prow$ remains incomplete
  (\cref{prop:robust-awareness}), and that incompleteness is subjectively
  optimal rather than merely mechanically absorbing.
  \item Firm $2$'s record contains the same four labels; its committed column
  policy disagrees at $x_{01}$ and $x_{10}$, so $\lvert L_2\rvert=2$ and
  $B_2=1$---twice the discovery-stage leverage.  Firm $2$ experiences
  dissonance, elects inquiry whenever $\theta_2\geq\gamma_2$, and, conditional
  on success, proceeds by \cref{prop:path-dependence} toward complete
  awareness.
\end{enumerate}
Differentiated defeat thus concentrates the incentive to inquire in the
defeated firm: the same episodes that transfer value to the correctly
formulated rival supply the loser with dissonance, revealed labels, and
leverage.  Competitive failure is, in this precise sense, an engine of
awareness growth, and sustained success rationally extinguishes the question.
\end{corollary}

\begin{proof}
The labels revealed across the two batches are $x_{00}\mapsto0$,
$x_{11}\mapsto1$ (discovery) and $x_{01}\mapsto0$, $x_{10}\mapsto1$
(validation).  Firm $1$'s policy $(0,0,1,1)$ agrees with every label, so
$L_1=\emptysetset$; absence of dissonance follows because
$m^r\in\widehat M(\Prow,D(h))$.  Firm $2$'s policy $(0,1,0,1)$ disagrees at
exactly the two off-diagonal states, so $\lvert L_2\rvert=2$.  Its dissonance
follows because under $\Pcol$ the label $x_{01}\mapsto0$ conflicts with
$x_{11}\mapsto1$ in one cell and $x_{10}\mapsto1$ conflicts with
$x_{00}\mapsto0$ in the other, emptying the active set.  The election claim
applies \cref{prop:participation}.
\end{proof}

\begin{remark}[Relation to deliberation-cost models]
\label{rem:cplus-relation}
The participation rule instantiates, in partition form, the deliberation logic
of the C-family agents of \cite{RyallAction2026}: refinement effort is
expended only when its represented value exceeds its cost.  There, the value
of investigating an unexplored branch is computed from objective payoff bounds
attached to unseen continuations.  The present treatment replaces those bounds
with the firm's own revealed record, in keeping with the discipline that
nothing behind the awareness frontier may carry objective payoff information
into pre-insight deliberation.
\end{remark}

\begin{remark}[An asymmetric partial-correction benchmark]
Suppose instead that firm $2$ remains exogenously fixed at architecture $1$ and
firm $1$ initially follows the false column theory.  A differentiated loss at
$x_{10}$ can lead firm $1$ to the partition
\[
\big\{\{x_{00}\},\{x_{10}\},\{x_{01},x_{11}\}\big\}
\]
and theory $(0,1,1,1)$.  Firm $1$ then wins at $x_{00}$ and pools elsewhere,
earning $5/8$ under equal state weights.  Its theory remains wrong at $x_{01}$,
where pooling conceals the error.  This calculation is useful as an
illustration, but it is not an equilibrium theorem unless firm $2$'s continued
fixed behavior and failure of inquiry are endogenized.
\end{remark}

\section{An outside option: value destruction and the knife-edge of
self-confirmation}
\label{sec:outside-option}

Under \cref{eq:objective-payoff}, total value is one in every episode: pooled
firms split the market even when both offer the objectively inferior
architecture, so misservice reallocates nothing and destroys nothing.  Three
results lean on that feature: the self-confirming trap
(\cref{thm:self-confirming-trap}), its immunity to cheap inquiry
(\cref{cor:trap-immunity}), and the permanent warrant gap
(\cref{cor:permanent-gap}).  This section parameterizes the demand side and
shows that all three are knife-edge properties of perfectly forgiving demand.

\begin{definition}[Outside option]
\label{def:outside-option}
For $\beta\in[0,1]$, replace \cref{eq:objective-payoff} with
\begin{equation}
u_i^\beta(a_i,a_j,x)=
\begin{cases}
\frac12, & a_i=a_j=d(x),\\[2pt]
\frac{\beta}{2}, & a_i=a_j\neq d(x),\\[2pt]
1, & a_i\neq a_j\text{ and }a_i=d(x),\\[2pt]
0, & a_i\neq a_j\text{ and }a_j=d(x).
\end{cases}
\label{eq:beta-payoff}
\end{equation}
Payoffs sum to the value created: one whenever some firm offers $d(x)$, and
$\beta$ when both firms offer the other architecture, in which case the
customer withholds $1-\beta$ of the unit value.  The subjective payoff
predicted by theory $m$ extends accordingly:
\begin{equation}
u_i^{m,\beta}(a_i,a_j,x)=
\begin{cases}
\frac12, & a_i=a_j\text{ and }m(x)\in\{a_i,n\},\\[2pt]
\frac{\beta}{2}, & a_i=a_j\text{ and }m(x)\notin\{a_i,n\},\\[2pt]
1, & a_i\neq a_j\text{ and }m(x)=a_i,\\[2pt]
0, & a_i\neq a_j\text{ and }m(x)=a_j,\\[2pt]
\frac12, & a_i\neq a_j\text{ and }m(x)=n.
\end{cases}
\label{eq:beta-subjective}
\end{equation}
The payoff rule, including $\beta$, is known to both firms; only the causal
law is not.  Every preceding section is the special case $\beta=1$.
\end{definition}

\begin{lemma}[Dominance under an outside option]
\label{lem:beta-dominance}
At an awareness cell $C$, the subjective expected-payoff differences between
actions $0$ and $1$ are
\[
v_i(0\mid C)-v_i(1\mid C)
=
\begin{cases}
\tfrac12\,p_i^0(C)-\bigl(1-\tfrac{\beta}{2}\bigr)p_i^1(C),
& \text{rival plays }0,\\[4pt]
\bigl(1-\tfrac{\beta}{2}\bigr)p_i^0(C)-\tfrac12\,p_i^1(C),
& \text{rival plays }1,
\end{cases}
\]
with mixtures handled by linearity.  Consequently:
\begin{enumerate}[label=(\roman*)]
  \item a firm certain of an architecture label finds the labeled action
  strictly dominant on that cell for every $\beta\in[0,1]$;
  \item action $0$ is strictly dominant on $C$ if and only if
  $p_i^0(C)>(2-\beta)\,p_i^1(C)$, and symmetrically for action $1$; hence
  belief-induced differentiation (\cref{prop:initial-differentiation}) extends
  to $\beta<1$ provided the prior asymmetry in
  \cref{ass:heterogeneous-priors} is strengthened to this threshold.
\end{enumerate}
\end{lemma}

\begin{proof}
Against rival action $0$:
$v_i(0\mid C)=\tfrac12[p_i^0(C)+p_i^n(C)]+\tfrac{\beta}{2}p_i^1(C)$ and
$v_i(1\mid C)=p_i^1(C)+\tfrac12 p_i^n(C)$.  Against rival action $1$:
$v_i(0\mid C)=p_i^0(C)+\tfrac12 p_i^n(C)$ and
$v_i(1\mid C)=\tfrac12[p_i^1(C)+p_i^n(C)]+\tfrac{\beta}{2}p_i^0(C)$.
Subtracting gives the displayed differences, which reduce to
\cref{eq:dominance-difference} at $\beta=1$.  For (i), $p_i^a(C)=1$ makes both
differences equal to $\tfrac12$ and $1-\tfrac{\beta}{2}$, both strictly
positive.  For (ii), both differences are strictly positive exactly when
$\tfrac12 p_i^0(C)>(1-\tfrac{\beta}{2})p_i^1(C)$.
\end{proof}

When $\beta<1$ and neither threshold in part (ii) is met, no constant action
is dominant and best responses acquire an anti-coordination structure:
differentiating guarantees that the customer is served and value one is
created, so demand-side harshness pushes balanced-belief play toward the
diagnostic regime.  A full equilibrium analysis of that game is deferred.

\begin{lemma}[Demand-side revelation]
\label{lem:demand-revelation}
Let $\beta<1$ and consider a pooled episode at state $x$ with common action
$a$.  The payoff level reveals the objective label: payoff one-half implies
$d(x)=a$ and excludes the theory value $1-a$ at $x$, while payoff $\beta/2$
implies $d(x)=1-a$ and excludes both the value $a$ and neutrality.
Differentiated episodes reveal labels as in \cref{lem:revelation}.  At
$\beta=1$, pooled episodes reveal nothing.
\end{lemma}

\begin{proof}
Objectively, \cref{eq:beta-payoff} gives pooled payoff one-half if $a=d(x)$
and $\beta/2$ otherwise, and the two are distinguishable exactly when
$\beta\neq1$.  Subjectively, \cref{eq:beta-subjective} predicts one-half under
$m(x)\in\{a,n\}$ and $\beta/2$ under $m(x)=1-a$; comparing predicted with
realized payoffs yields the stated exclusions.
\end{proof}

\begin{remark}[Leverage and participation under an outside option]
\label{rem:leverage-beta}
For $\beta<1$, \cref{def:disagreement-set} extends by counting as labeled any
pooled episode whose payoff level pins the objective label
(\cref{lem:demand-revelation}).  The unit gain of implementing a disagreeing
label is no longer conjecture-free: depending on the rival's action at the
recurrence, it lies in $[\tfrac12,\,1-\tfrac{\beta}{2}]$, so
$B_i(h)=\lvert L_i(h)\rvert/2$ is a conjecture-free \emph{lower bound} on the
pinned gain rather than its exact value, and the participation rule of
\cref{prop:participation} errs only toward abstention.  The results of
\cref{sec:inquiry-economics} hold as stated at $\beta=1$; for $\beta<1$ the
trap-immunity corollary loses its premise because, as shown next, the trap
itself fails.
\end{remark}

\begin{proposition}[No payoff-relevant self-confirmation for $\beta<1$]
\label{prop:no-trap}
Let $\beta<1$.  Consider awareness partitions, degenerate beliefs (each firm
certain of a single expressible theory), and committed policies recommitted
each generation, satisfying conditions (i) and (ii) of
\cref{def:awareness-trap}---subjective optimality and on-path
consistency---and suppose every customer state recurs.  Then every state is
served by its objectively preferred architecture, so value creation is
maximal in every generation, and every architecture label is objectively
correct: $m_i(x)\in\{0,1\}$ implies $m_i(x)=d(x)$.  Residual falsity is
confined to neutral mislabeling---states labeled $n$ whose realized service is
nonetheless correct---which affects neither value creation nor value capture.
\end{proposition}

\begin{proof}
Fix a recurring state $x$.  First, the actions must pool at $x$.  Suppose they
differ; the losing firm observes payoff zero.  If the loser's cell carries an
architecture label, \cref{lem:beta-dominance}(i) and subjective optimality
require its action to equal that label, so its predicted payoff is one, and if
its cell is labeled $n$ its predicted payoff is one-half; either way the
observation contradicts consistency.  Second, suppose the pooled payoff is
$\beta/2$.  By \cref{lem:demand-revelation}, consistency requires
$m_i(x)=1-a$ for both firms; \cref{lem:beta-dominance}(i) then prescribes
action $1-a$, contradicting the pooled action $a$.  Hence the pooled payoff is
one-half, so $d(x)=a$: the state is correctly served, and consistency admits
only $m_i(x)\in\{a,n\}$, so any architecture label equals $d(x)$.
\end{proof}

\begin{corollary}[Knife-edge of the epistemic pathologies]
\label{cor:knife-edge}
The column profile of \cref{thm:self-confirming-trap} is a self-confirming
awareness trap if and only if $\beta=1$.  For $\beta<1$, in the first
generation containing an off-diagonal state:
\begin{enumerate}[label=(\roman*)]
  \item both firms receive $\beta/2$ where their shared theory predicts
  one-half; the pinned label conflicts with a previously revealed diagonal
  label inside a cell of $\Pcol$, the active set empties, and
  \cref{ass:dissonance-trigger} governs renewed inquiry;
  \item the pinned label disagrees with the committed column policy, so the
  extended disagreement set is nonempty, $B_i>0$, and the zero-leverage lock
  of \cref{cor:trap-immunity} fails;
  \item the observed payoff has likelihood one under $d^r$ and zero under
  $d^c$, so the reflective posterior is degenerate and the permanent warrant
  gap of \cref{cor:permanent-gap} fails.
\end{enumerate}
The destroyed value $1-\beta$ at each misserved episode is itself the
corrective signal: what demand withholds from the firms, their epistemics
receives as evidence.
\end{corollary}

\begin{proof}
The $\beta=1$ direction is \cref{thm:self-confirming-trap}.  Let $\beta<1$
and let the off-diagonal state be $x_{01}$ (the case $x_{10}$ is symmetric).
The column policies pool on action $1$ there while $d^r(x_{01})=0$, so the
realized payoff is $\beta/2$, whereas $m^c(x_{01})=1$ predicts one-half:
consistency fails, so the profile is not a trap.  By
\cref{lem:demand-revelation} the label $x_{01}\mapsto0$ is pinned; it
conflicts with the discovery label $x_{11}\mapsto1$ in the cell
$\{x_{01},x_{11}\}$ of $\Pcol$, so the active set empties
(\cref{prop:conflict-characterization}).  The committed policy plays $1$ at
$x_{01}$, so the episode enters the extended disagreement set and
$B_i\geq\tfrac12$.  Finally, under $d^c$ the pooled action at $x_{01}$ is
correct and the predicted payoff is one-half, so the realized $\beta/2$ has
likelihood zero, while under $d^r$ it has likelihood one; the reflective
posterior concentrates on $d^r$.
\end{proof}

\begin{remark}[Two sources of diagnosticity]
\label{rem:two-sources}
At $\beta=1$, evidence arises only from strategic differentiation: theory
heterogeneity across firms is the sole corrective force, and consensus on a
false formulation is absorbing
(\cref{thm:self-confirming-trap,cor:trap-immunity,cor:permanent-gap}).  For
$\beta<1$, the demand side supplies evidence precisely when the strategy side
pools, and \cref{prop:no-trap} shows that payoff-relevant collective error
cannot persist.  The comparative interpretation is direct: markets with
forgiving demand---captive customers, absent substitutes, buyers who do not
bear the cost of misservice---require competitive variety in beliefs to
correct false collective theories, whereas markets with strong outside
options discipline theories even under consensus.  For the general dynamic
model, the parameter locates two learning regimes: $\beta<1$ supplies an
evidence stream independent of strategic heterogeneity, while $\beta=1$ makes
all learning contingent on it.
\end{remark}

\section{Cognitional interpretation}

The model's relationship to Lonergan's account of inquiry is interpretive, not
algorithmic.  Lonergan does not supply a partition lattice, a coordinate-
complexity criterion, or a Bayesian rule.  His distinctions instead discipline
which formal objects should not be conflated.  The relevant distinctions among
insight, formulation, and reflective judgment are developed in
Lonergan~\cite[chs.~1 and 10]{Lonergan1992}; the question--insight--formulation--
validation sequence is also developed in Ryall and
Wilkins~\cite{RyallWilkins2018}.

\begin{center}
\renewcommand{\arraystretch}{1.2}
\begin{tabularx}{\textwidth}{@{}>{\raggedright\arraybackslash}p{0.22\textwidth}>{\raggedright\arraybackslash}p{0.34\textwidth}X@{}}
\toprule
\textbf{Cognitional operation} & \textbf{Formal counterpart} &
\textbf{Qualification}\\
\midrule
Experience & Customer episodes, actions, and outcomes & Particular experience
does not itself supply a causal predicate.\\
Question for intelligence & In this baseline, causal dissonance makes salient
the question of what explains the reversal & The question specifies the defect
and requirements of an answer, not a known answer menu.\\
Heuristic anticipation & Restoration, minimality, and projectibility & These
are the model's analogue of heuristic anticipation: they constrain what an
answer must accomplish without providing its semantic content.\\
Direct insight & No direct mathematical counterpart & The model represents
only the ensuing representational and formulational products, not the conscious
occurrence.\\
Representation production & New partition $\mathcal P'$ & A previously
unavailable distinction becomes usable in theory and action.\\
Formulation & Pair $(\mathcal P',m')$ & A distinction alone is not a causal
explanation; a stable relation must be articulated on it.\\
Question for reflection & Whether $m'$ holds beyond the generating cases &
Retrospective fit is not independent confirmation.\\
Judgment & Assessment of whether the formulation is adequately warranted &
Fresh diagnostic evidence or a selection-adjusted likelihood supplies the
model's implementation.\\
Decision & Policy measurable with respect to $\mathcal P'$ & Strategic action
determines which future experience will be diagnostic; provisional action may
occur before judgment is complete.\\
\bottomrule
\end{tabularx}
\end{center}

In this interpretation, the formal sequence associated with a successful
insight is not exhausted by partition production: the available distinction
must subsequently be expressed in a projectible causal formulation.
Restoration is weaker than truth, so the candidate formulation
remains subject to reflective judgment.  Agent-specific prior understanding,
analogy, habits of comparison, and organizational routines can eventually be
modeled as differences in heuristic accessibility or in the complexity
ordering, but they must not become a subjective menu over semantically named
answers before those answers are thinkable.

In Lonergan's account the desire to know is unrestricted, but its deliberate
prosecution by an economic agent is not free.  The participation layer of
\cref{sec:inquiry-economics} separates the spontaneous emergence of a question
(dissonance makes the defect salient without any choice) from the deliberate
allocation of effort to pursue it (the election $\varepsilon_i$).  The former
is a cognitive event; the latter is an economic act, and the leverage
statistic shows that it can be priced without violating the stricture that
pre-insight deliberation must not quantify over unformulated answers.

Causal dissonance should not be identified with inverse insight.  A more
appropriate later analogue of inverse insight would be a transition in which
the agent grasps that another deterministic customer partition is the wrong
kind of explanation---for example, because the relevant law is stochastic,
strategic, or located at a different causal level.  That transition lies beyond
the deterministic four-state model.

\section{Results, limitations, and the next theoretical problem}

The formal results have different roles.  The following are necessary
scaffolding:
\begin{enumerate}[label=(\roman*)]
  \item private causal beliefs can induce strategically differentiated product
  commitments;
  \item differentiated actions reveal causal labels while pooled actions do
  not;
  \item fixed-space Bayesian conditioning cannot produce a new awareness
  partition;
  \item conflict-free cells characterize restorative refinements; and
  \item the four-state discovery history has exactly four minimal restorative
  refinements.
\end{enumerate}

The potentially substantive mechanisms are sharper:
\begin{enumerate}[label=(\roman*)]
  \item \emph{Fallibility versus exhaustiveness.}  Observationally equivalent
  worlds preclude a uniformly correct truth-blind formulation, while a
  representation robustly sufficient for both row and column worlds must be
  completely fine.
  \item \emph{Representation-relative certainty.}  A false causal formulation
  can receive probability one under retrospective renormalization inside its
  newly produced model space.
  \item \emph{Path-dependent awareness.}  Firms can converge to the same true
  causal theory and policy while retaining different awareness structures.
  \item \emph{Endogenous diagnosticity.}  Theories determine actions, actions
  determine which causal labels experience reveals, and revealed labels
  determine whether theories survive.
  \item \emph{Self-confirming causal homogeneity.}  Shared false formulations
  can align strategies and suppress their own falsification, whereas causal
  heterogeneity generates both learning and asymmetric value capture.
  \item \emph{Delabeled question value.}  The subjective value of a question
  is computable from retained episodes alone, is invariant across every
  representation that could answer it, and tracks action separation rather
  than residual model multiplicity.  Consequently the awareness trap is robust
  to cheap inquiry, sustained competitive success rationally extinguishes
  inquiry, and differentiated defeat makes the losing firm the natural
  inquirer.
  \item \emph{Unwarranted certainty and its invisibility.}  Retrospective fit
  is guaranteed by restoration and hence uninformative; representation-
  relative certainty exceeds selection-adjusted warrant by a gap that pooled
  play preserves forever; and no firm short of complete awareness can express
  the comparison that measures its own gap.  Question value and warrant
  improvement share a currency: both are created by differentiated action and
  destroyed by pooling.
  \item \emph{Knife-edge pathologies.}  The self-confirming trap, the
  permanent warrant gap, and the zero-leverage lock all require demand to
  absorb misservice completely.  With any outside option, consistency and
  subjective optimality force correct service of every recurring state; false
  theories persist only as payoff-irrelevant neutral mislabels, and the
  destroyed value is itself the corrective signal.  Diagnosticity has two
  sources---strategic heterogeneity and demand-side harshness---and collective
  error requires the absence of both.
\end{enumerate}

The example does not yet supply a complete strategy-field contribution.  In
particular, it endogenizes only the binary participation margin of inquiry,
not its intensity or duration; and it does not endogenize the production
of a bridge prior, the accumulation of agent-specific heuristic accessibility,
or the organizational and competitive conditions that make alternative
inquiry technologies valuable.  \Cref{sec:outside-option} connects false
awareness to value creation through the outside option, but the demand
response itself remains exogenous; endogenizing $\beta$ through prices,
customer search, or exit dynamics---and characterizing equilibrium in the
anti-coordination region where no constant action is dominant---remains open.

The next theoretical problem is therefore not simply to enlarge $X$.  It is to
endogenize the representation-producing technology and its strategic
environment while preserving the four-state model's central discipline:
experience may direct inquiry, but neither inquiry nor Bayes may be allowed to
select the objective truth by assumption.

\section{Conclusion}

The $2\times2$ environment isolates a formal boundary between learning within a
theory space and producing a new theory space.  Strategic differentiation
creates a causal conflict that no initially expressible theory can resolve.
Restoration and parsimony narrow the possible responses from fifteen partitions
to two simple factor-based candidates, but observational equivalence prevents
the experience from selecting the true one.  A non-exhaustive inquiry
technology must therefore be fallible.  Once a formulation is produced,
fixed-space Bayesian assessment can resume, but representation-relative
certainty must be separated from judgment based on adequate fresh or
selection-adjusted evidence.  Finally, strategic action feeds back into the epistemic process:
shared false formulations may become self-confirming by eliminating diagnostic
variation, whereas theory heterogeneity generates both competitive rents and
the experience required for correction.  Both sides of that contrast
presuppose perfectly forgiving demand; with any outside option, the customer
side supplies the discipline that pooled strategies suppress.

\begin{thebibliography}{9}

\bibitem{Lonergan1992}
Lonergan, Bernard J. F. 1992.
\newblock \emph{Insight: A Study of Human Understanding}.
\newblock Collected Works of Bernard Lonergan, vol.~3, edited by Frederick E.
Crowe and Robert M. Doran. Toronto: University of Toronto Press.

\bibitem{RyallAction2026}
Ryall, M. D. 2026.
\newblock ``Mathematical Analysis of Agents $C$, $C^{+}$, and
$C^{\mathrm{Plan}}$.''
\newblock Unpublished manuscript, Appendix A.

\bibitem{RyallWilkins2018}
Ryall, M. D., and J. Wilkins. 2018.
\newblock ``Insight and Social Being.''
\newblock Preliminary manuscript.

\end{thebibliography}

\end{document}
